% =====================================================================
%  IEEE-STYLE CONFERENCE REPORT TEMPLATE  (self-contained)
%  MLE-AI-201 Machine Learning -- Final Project
%
%  This template MIMICS the IEEE two-column conference look using only
%  standard LaTeX packages. It does NOT depend on IEEEtran.cls, so it
%  compiles anywhere (tectonic, pdflatex, Overleaf) with nothing extra.
%
%  Build:   tectonic report.tex      (recommended)
%     or:   pdflatex report.tex ; pdflatex report.tex
% =====================================================================
\documentclass[10pt,twocolumn,a4paper]{article}

% ---- Geometry: IEEE-like margins / two narrow columns ----------------
\usepackage[a4paper,top=0.75in,bottom=1.0in,left=0.68in,right=0.68in,columnsep=0.25in]{geometry}

% ---- Times-like font (standard, ships with every TeX distro) ---------
\usepackage{mathptmx}          % Times text + math
\usepackage[T1]{fontenc}
\usepackage[utf8]{inputenc}    % allows accented / Azerbaijani chars

% ---- Common, dependency-free packages --------------------------------
\usepackage{graphicx}
\usepackage{amsmath,amssymb}
\usepackage{booktabs}
\usepackage{array}
\usepackage{titlesec}
\usepackage{caption}
\usepackage{cite}
\usepackage[hidelinks]{hyperref}
\usepackage{url}

% ---- IEEE-style section headings -------------------------------------
\renewcommand{\thesection}{\Roman{section}}
\renewcommand{\thesubsection}{\Alph{subsection}}
\titleformat{\section}[block]{\centering\normalsize\scshape}{\thesection.}{0.5em}{}
\titleformat{\subsection}[block]{\normalsize\itshape}{\thesubsection.}{0.5em}{}
\titlespacing*{\section}{0pt}{1.4ex plus 1ex minus .2ex}{1ex}
\titlespacing*{\subsection}{0pt}{1.0ex plus 1ex minus .2ex}{0.6ex}

% ---- Abstract / Index-Terms in the IEEE boldface-dash style ----------
\newcommand{\abstractkw}[1]{\noindent\textbf{\textit{#1}}}

\captionsetup{font=small,labelsep=period}
\setlength{\parindent}{1em}

% =====================================================================
\begin{document}

% ------------------------- TITLE BLOCK --------------------------------
\twocolumn[
  \begin{@twocolumnfalse}
  \begin{center}
    {\LARGE\bfseries Real-Estate Price Prediction:\\
    From-Scratch Trees and Pegasos SVMs\par}
    \vspace{1.2em}
    % Authors: Surname, Given name --- alphabetical by surname.
    {\large
      Babayev, Mahhamadali\quad Hajiyev, Ilham\quad
      Javadzade, Mahammad\quad Isgandarli, Fariz\par}
    \vspace{0.5em}
    {\itshape MLE-AI-201 Machine Learning, AI Academy --- Cohort II\par}
    \vspace{1.4em}
  \end{center}
  \begin{quote}
    \small
    \abstractkw{Abstract---}% IEEE uses an em-dash after "Abstract"
    We study Azerbaijani real-estate listings from bina.az through two
    supervised tasks: predicting price in AZN and classifying listings as
    standard or premium relative to the training median. The pipeline
    cleans mixed-type fields, removes target-derived leakage, encodes
    categoricals, and fits imputation and scaling using training data only.
    We implement a NumPy decision tree and a Pegasos soft-margin SVM, with
    Random Fourier Features (RFF) as a nonlinear extension, and compare
    them with scikit-learn baselines. On the held-out test split, the
    scratch tree obtains RMSE 380,596 AZN and $R^2=0.688$ for regression;
    the scratch linear Pegasos model obtains F1 0.826 and ROC-AUC 0.911.
    The scikit-learn tree gives the highest classification F1 (0.930) and
    ROC-AUC (0.967). Linear Pegasos remains competitive with linear
    baselines, whereas RFF is sensitive to bandwidth and finite feature
    dimension. Code is available at
    \href{https://github.com/Haji708/ml_final_project}{\texttt{github.com/Haji708/ml\_final\_project}}.

    \vspace{0.6em}
    \abstractkw{Index Terms---}decision tree, support-vector machine,
    Pegasos, regression, classification, real-estate price prediction.
  \end{quote}
  \vspace{1.0em}
  \end{@twocolumnfalse}
]

% =====================================================================
\section{Introduction}
Online property markets contain heterogeneous numeric, categorical, and
geospatial signals. We frame the problem as Task A, price regression, and
Task B, binary price-tier classification. Our contributions are a shared
leakage-aware preprocessing pipeline, NumPy implementations of a decision
tree and Pegasos SVM, an RFF extension, and matched scikit-learn benchmarks.

\section{Related Work / Background}
The tree recursively chooses numeric thresholds that maximize weighted
impurity reduction. Gini impurity and entropy are used for classification,
while variance/MSE reduction is used for regression. Leaves predict the
majority class and class frequencies, or the mean target. The SVM minimizes
the regularized hinge-loss primal with stochastic sub-gradient updates.
Pegasos avoids a quadratic-programming solver and uses
$\eta_t=1/(\lambda t)$~\cite{pegasos}. RFF approximates an RBF feature map
before applying the same linear optimizer.

\section{Data and Preprocessing}
The downloaded bina.az snapshot contains 100,775 raw listings. We enforce
AZN prices, coerce numeric strings, parse mixed units, and remove 34
records outside the fixed eligibility range
$1{,}000\leq\text{price}\leq20{,}000{,}000$ AZN. Exact duplicates and
duplicate identifiers were absent. Missing numeric values use
training-fitted medians and missingness indicators; categorical values are
fitted from the training partition. The resulting modelling dataset has
100,705 rows and 178 encoded features.

Task A predicts price and uses RMSE, MAE, and $R^2$. Task B defines
premium as price above the training median of 218,000 AZN and standard
otherwise; it uses precision, recall, F1, and ROC-AUC. The deterministic
seed-42 split is 70,493/15,106/15,106 train/validation/test rows, sampled
within price-rank strata. This balances price ranks but is not exact binary
tier stratification. The leakage columns \texttt{unit\_price} and
\texttt{total\_price} are removed because they reveal or derive the target.
IDs, URLs, owner names, addresses, and free text are also excluded.

\begin{table}[t]
  \centering
  \caption{Dataset summary after cleaning.}
  \label{tab:data}
  \begin{tabular}{@{}lr@{}}
    \toprule
    Property & Value \\
    \midrule
    Rows (after cleaning)       & 100,705 \\
    Features used               & 178 encoded features \\
    Train / Val / Test split    & 70,493 / 15,106 / 15,106 \\
    Training tier balance      & 50.14\% / 49.86\% \\
    \bottomrule
  \end{tabular}
\end{table}

\section{Method}
\subsection{Decision tree (from scratch)}
For each candidate feature, values are sorted and thresholds are tested
between distinct adjacent values. The selected split maximizes parent
impurity minus the sample-weighted child impurities. Growth is recursive
and stops at maximum depth, insufficient samples, invalid child sizes, or
non-positive impurity reduction. Classification leaves store class
probabilities and predict the majority class; regression leaves predict
the target mean. Validation experiments use depth and minimum-leaf
controls to limit overfitting, with Gini, entropy, or MSE selected
according to the task.

\subsection{Support-vector machine (from scratch)}
The objective and Pegasos update are:
\begin{equation}
  J(\mathbf{w},b)=\tfrac{\lambda}{2}\lVert\mathbf{w}\rVert^{2}
  +\tfrac{1}{n}\sum_{i=1}^{n}\max\!\big(0,\,1-y_i(\mathbf{w}^{\top}\mathbf{x}_i+b)\big).
\end{equation}
Labels are mapped to $\{-1,+1\}$. At step $t$, a sampled example with
margin $m_i=y_i(\mathbf{w}^{\top}\mathbf{x}_i+b)<1$ updates both weights
and bias; otherwise only regularization shrinkage is applied to
$\mathbf{w}$. The implementation records the objective during fitting.
All SVM inputs are standardized using training means and scales only.
The nonlinear extension uses RFF:
\[
\phi(\mathbf{x})=\sqrt{2/D}\cos(\mathbf{x}W+\mathbf{r}),
\]
where $W$ and $\mathbf{r}$ are generated once during fitting and reused
at prediction time. Support vectors satisfy
$y_i f(\mathbf{x}_i)\leq1+10^{-6}$, and the geometric margin width is
$2/\lVert\mathbf{w}\rVert$.

\section{Experimental Setup}
We report RMSE, MAE, and $R^2$ for Task A, and precision, recall, F1,
ROC-AUC, and confusion-matrix diagnostics for Task B. Hyperparameters are
selected on validation data only. The linear Pegasos sweep tests
$\lambda\in\{10^{-5},10^{-4},10^{-3}\}$; RFF trials vary bandwidth and
component count. The selected linear model uses $\lambda=10^{-3}$,
50,000 updates, and the fixed seed is 42. Baselines include
normal-equation Ridge, scikit-learn Ridge, DecisionTreeRegressor/
Classifier, SGDClassifier with hinge loss, and LinearSVC.

\section{Results}
Lead with a results table; back every claim with a number or figure.

\begin{table}[t]
  \centering
  \caption{Headline results on the held-out test set.}
  \label{tab:results}
  \begin{tabular}{@{}lcc@{}}
    \toprule
    Model & Task A (RMSE) & Task B (F1) \\
    \midrule
    Ours: decision tree     & 380,596 & 0.941 \\
    Ours: SVM (Pegasos)     & n/a & 0.826 \\
    scikit-learn baseline   & 401,341 & 0.930 \\
    \bottomrule
  \end{tabular}
\end{table}

\begin{figure*}[!t]
  \centering
  \includegraphics[width=0.47\textwidth]{figures/regression_validation.png}
  \hfill
  \includegraphics[width=0.47\textwidth]{figures/classification_validation.png}
  \par\vspace{0.5em}
  \includegraphics[width=0.47\textwidth]{figures/svm_convergence.png}
  \hfill
  \includegraphics[width=0.47\textwidth]{figures/error_analysis.png}
  \caption{Validation performance, Pegasos objective convergence, and test
  error analysis. Regression RMSE is lower-is-better; classification F1 is
  higher-is-better.}
  \label{fig:main}
\end{figure*}

\section{Discussion and Limitations}
The scratch tree gives the best reported regression test performance
(RMSE 380,596 AZN, MAE 71,856 AZN, $R^2=0.688$), improving on the
scikit-learn tree's RMSE 401,341 AZN. For classification, the
scikit-learn tree is strongest overall (F1 0.930, ROC-AUC 0.967), while
linear Pegasos reaches F1 0.826 and ROC-AUC 0.911. Its result is close
to SGDClassifier (F1 0.838) and LinearSVC (F1 0.835), supporting the
correctness of the from-scratch linear solver.

RFF is a valid nonlinear extension but is sensitive to bandwidth, random
approximation error, and component count. On this snapshot its test F1
is 0.832 and ROC-AUC is 0.908, so the finite approximation did not
improve on the linear boundary. Large residuals are expected for luxury
listings and sparse locations; a single global model cannot fully
represent local market effects. The fixed split is not cross-validated,
so these values quantify this snapshot rather than uncertainty across
folds. Future work should add repeated or temporal validation, richer
geospatial features, and systematic RFF bandwidth tuning.

\section{Conclusion}
The project demonstrates a complete leakage-aware workflow for bina.az
price prediction and price-tier classification. From-scratch Pegasos
matches the expected hinge-loss formulation and remains competitive with
linear scikit-learn baselines, while the custom tree captures stronger
nonlinear structure in this dataset. The next priority is repeated or
time-based validation and fuller analysis of regional and luxury-listing
errors.

\section*{Tools and Acknowledgements}
The team used NumPy, pandas, matplotlib, and scikit-learn for data
processing, modelling, evaluation, and plotting. An AI coding assistant
was used for code navigation, debugging suggestions, LaTeX editing, and
test-oriented review; all generated changes were inspected and validated
by the team. The raw bina.az data were downloaded from the Kaggle
\texttt{binaaz-sale-project} dataset by Sehriyar Memmedli and are not
redistributed in this repository.

% =====================================================================
\begin{thebibliography}{9}
\bibitem{pegasos}
S.~Shalev-Shwartz, Y.~Singer, N.~Srebro, and A.~Cotter,
``Pegasos: primal estimated sub-gradient solver for SVM,''
\emph{Mathematical Programming}, vol.~127, no.~1, pp.~3--30, 2011.
\end{thebibliography}

\end{document}
